\documentclass[aspectratio=169,10pt]{beamer}
% ---------------------------------------------------------------
% R4.02 Qualite de developpement -- preambule commun des diapositives
% Charte reprise de R3.07 (metropolis, dark tech)
% Palette IUT partagee (violet du logo + accent or)
% Utilisation :
%   \documentclass[aspectratio=169,10pt]{beamer}
%   % ---------------------------------------------------------------
% R4.02 Qualite de developpement -- preambule commun des diapositives
% Charte reprise de R3.07 (metropolis, dark tech)
% Palette IUT partagee (violet du logo + accent or)
% Utilisation :
%   \documentclass[aspectratio=169,10pt]{beamer}
%   % ---------------------------------------------------------------
% R4.02 Qualite de developpement -- preambule commun des diapositives
% Charte reprise de R3.07 (metropolis, dark tech)
% Palette IUT partagee (violet du logo + accent or)
% Utilisation :
%   \documentclass[aspectratio=169,10pt]{beamer}
%   \input{../../latex/common/r402-slides-preamble.tex}
%   \graphicspath{{../../latex/common/assets/}{img/}{figs/}}
% ---------------------------------------------------------------
\usetheme{metropolis}
\usepackage{appendixnumberbeamer}

\usepackage[T1]{fontenc}
\usepackage[utf8]{inputenc}
\usepackage[french]{babel}
\usepackage[scale=0.94]{tgheros}
\renewcommand{\familydefault}{\sfdefault}
\usepackage{amsmath,amssymb,amsfonts}
\usepackage{graphicx}
\usepackage{booktabs}
\usepackage{multirow}
\usepackage{tikz}
\usetikzlibrary{shapes.geometric, shapes, arrows, arrows.meta, positioning, calc, fit}
\usepackage{xcolor}
\usepackage{listings}
\usepackage[skins,listings]{tcolorbox}
\usepackage{qrcode}
\usepackage{eurosym}
\usepackage{hyperref}

% ---------------------------------------------------------------
% PALETTE (dark-mode tech: near-black canvas, violet (logo IUT) + gold accents)
% ---------------------------------------------------------------
\definecolor{qabg}{HTML}{0E1016}      % canvas background
\definecolor{qapanel}{HTML}{181B24}   % card / block panel
\definecolor{qapanelalt}{HTML}{20242F}% secondary panel tint
\definecolor{qaline}{HTML}{2B303D}    % hairline rules / table rules
\definecolor{qatext}{HTML}{ECEEF2}    % primary text (off-white)
\definecolor{qamuted}{HTML}{8B93A6}   % secondary / caption text
\definecolor{qaaccent}{HTML}{D93ED7}  % primary accent -- violet du logo IUT, eclairci pour fond sombre
\definecolor{qaaccent2}{HTML}{F2B84B} % secondary accent -- or/ambre
\definecolor{qacodekw}{HTML}{C792EA}  % code : mots-cles
\definecolor{qacodestr}{HTML}{C3E88D} % code : chaines
\definecolor{qacodecom}{HTML}{7F8AA3} % code : commentaires
\definecolor{qacodeann}{HTML}{F2B84B} % code : annotations / javadoc

\setbeamercolor{background canvas}{bg=qabg}
\setbeamercolor{normal text}{fg=qatext}
\setbeamercolor{palette primary}{bg=qabg, fg=qamuted}
\setbeamercolor{frametitle}{bg=qabg, fg=qatext}
\setbeamercolor{progress bar}{fg=qaaccent, bg=qaline}
\setbeamercolor{alerted text}{fg=qaaccent2}
\setbeamercolor{block title}{bg=qapanel, fg=qaaccent}
\setbeamercolor{block body}{bg=qapanel, fg=qatext}
\setbeamercolor{block title alerted}{bg=qapanel, fg=qaaccent2}
\setbeamercolor{block body alerted}{bg=qapanel, fg=qatext}
\setbeamercolor{block title example}{bg=qapanelalt, fg=qaaccent}
\setbeamercolor{block body example}{bg=qapanelalt, fg=qatext}
\setbeamercolor{example text}{fg=qaaccent}
\setbeamercolor{itemize item}{fg=qaaccent}
\setbeamercolor{itemize subitem}{fg=qaaccent2}
\setbeamercolor{enumerate item}{fg=qaaccent}
\setbeamercolor{enumerate subitem}{fg=qaaccent2}
\setbeamercolor{caption}{fg=qamuted}
\setbeamercolor{page number in head/foot}{fg=qamuted, bg=qabg}
\setbeamercolor{title}{fg=qatext}
\setbeamercolor{subtitle}{fg=qamuted}
\setbeamercolor{author}{fg=qatext}
\setbeamercolor{institute}{fg=qamuted}
\setbeamercolor{date}{fg=qamuted}
\setbeamercolor{section in toc}{fg=qaaccent}
\setbeamercolor{subsection in toc}{fg=qamuted}
% Sommaire de debut de section : section en cours en violet et gras, les autres en gris
\setbeamercolor{section in toc shaded}{fg=qamuted}
\setbeamercolor{subsection in toc shaded}{fg=qamuted}

\setbeamertemplate{blocks}[rounded][shadow=false]
\setbeamertemplate{itemize item}{\raisebox{0.15em}{\scriptsize\textcolor{qaaccent}{\ensuremath{\blacktriangleright}}}}
\setbeamertemplate{itemize subitem}{\raisebox{0.1em}{\tiny\textcolor{qaaccent2}{\ensuremath{\bullet}}}}

\setbeamerfont{frametitle}{series=\bfseries, size=\Large}
\setbeamerfont{title}{series=\bfseries, size=\Large}
\setbeamerfont{subtitle}{size=\normalsize}
\setbeamerfont{institute}{size=\scriptsize}
\setbeamerfont{date}{size=\scriptsize}
\setbeamerfont{author}{size=\small}

\metroset{sectionpage=none, subsectionpage=none, numbering=fraction, progressbar=frametitle}
\setbeamertemplate{section in toc}[sections numbered]

% ---------------------------------------------------------------
% LISTINGS -- theme sombre (lstlisting reste utilisable tel quel)
% \lstcode : police + couleur de base, a reutiliser dans les
%            surcharges locales (ex. basicstyle=\lstcode\scriptsize)
% ---------------------------------------------------------------
\newcommand{\lstcode}{\ttfamily\color{qatext}}

\lstdefinelanguage{TypeScript}{
  keywords={typeof, new, true, false, catch, function, return, null, switch, var,
    if, in, while, do, else, case, break, const, let, interface, class, extends,
    implements, enum, export, import, this, super, constructor, public, private,
    protected, static, async, await, throw, throws, try, finally, void},
  ndkeywords={boolean, any, number, string, int, long, double},
  sensitive=true,
  comment=[l]{//},
  morecomment=[s]{/*}{*/},
  morestring=[b]',
  morestring=[b]"
}

\lstset{
  language=Java,
  basicstyle=\lstcode\footnotesize,
  keywordstyle=\color{qacodekw}\bfseries,
  ndkeywordstyle=\color{qaaccent2}\bfseries,
  identifierstyle=\color{qatext},
  stringstyle=\color{qacodestr},
  commentstyle=\color{qacodecom}\itshape,
  morecomment=[s][\color{qacodeann}]{/**}{*/},
  backgroundcolor=\color{qapanelalt},
  rulecolor=\color{qaline},
  frame=single,
  framesep=3pt,
  xleftmargin=4pt,
  xrightmargin=4pt,
  keepspaces=true,
  columns=fullflexible,
  showstringspaces=false,
  breaklines=true,
  tabsize=2,
  escapeinside={(*@}{@*)},
  % accents UTF-8 dans le code (listings ne gere pas l'UTF-8 multi-octets)
  literate={é}{{\'e}}1 {è}{{\`e}}1 {ê}{{\^e}}1 {à}{{\`a}}1 {â}{{\^a}}1
           {ù}{{\`u}}1 {û}{{\^u}}1 {ô}{{\^o}}1 {î}{{\^i}}1 {ç}{{\c c}}1
           {É}{{\'E}}1 {À}{{\`A}}1 {€}{{\euro{}}}1 {…}{{\ldots}}1,
}

% ---------------------------------------------------------------
% CARD-STYLE BOXES
% ---------------------------------------------------------------
\tcbset{
  figcard/.style={colback=white, colframe=white, boxrule=0pt, arc=4pt,
    left=3pt, right=3pt, top=3pt, bottom=3pt, boxsep=0pt,
    halign=center, valign=center},
  datacard/.style={colback=qapanel, colframe=qaline, boxrule=0.4pt, arc=3pt,
    left=8pt, right=8pt, top=6pt, bottom=6pt, coltext=qatext,
    before upper={\color{qatext}}, halign=center},
  codebox/.style={colback=qapanelalt, colframe=qaline, boxrule=0.4pt, arc=3pt,
    left=8pt, right=8pt, top=5pt, bottom=5pt, listing only,
    listing options={basicstyle=\ttfamily\small\color{qaaccent},
      breaklines=true, showstringspaces=false, columns=fullflexible,
      keepspaces=true, frame=none, backgroundcolor={}}}
}
\newtcolorbox{figcard}[1][]{figcard,#1}
\newtcolorbox{datacard}{datacard}
\newtcblisting{codebox}{codebox}

% ---------------------------------------------------------------
% HANDY MACROS
% ---------------------------------------------------------------
% Mention d'elaboration (discrete : plus sombre que le texte secondaire)
\definecolor{qaaimention}{HTML}{4E5566}
\newcommand{\raimention}{{\tiny\color{qaaimention}Supports \'elabor\'es avec l'aide de Claude Code Opus 5.5}}
\newcommand{\hl}[1]{\alert{#1}}
\newcommand{\source}[1]{{\scriptsize\color{qamuted}#1}}

% Sommaire automatique en debut de section
\AtBeginSection[]{%
  \begin{frame}<beamer>{Plan du cours}
    \tableofcontents[currentsection]
  \end{frame}}

% Page de titre : \rtitleslide{numero du cours}{titre du cours}
% Les auteurs sont repris de \rauthors (modifiable avant \begin{document})
\newcommand{\rauthors}{Chouki Tibermacine}
\newcommand{\rtitleslide}[2]{%
\begin{frame}[plain,noframenumbering]
  \begin{center}
    \vspace{2mm}
    {\color{qaaccent}\rule{2.1cm}{1.2pt}}\\[6mm]
    {\Large\bfseries\color{qatext} R4.02}\\[1mm]
    {\Large\bfseries\color{qatext} Qualit\'e de d\'eveloppement}\\[5mm]
    {\large\color{qamuted} Cours #1}\\[1mm]
    {\large\color{qaaccent} #2}\\[6mm]
    \includegraphics[height=2.1cm]{iut_logo_transparent.png}

    \vspace{4mm}

    {\small\color{qatext} \rauthors}\\[1mm]
    {\scriptsize\color{qamuted} BUT 2 Informatique --- IUT de Vannes --- Universit\'e Bretagne Sud}\\[2mm]
    \raimention
  \end{center}
\end{frame}
}

% ---- Diapo de demonstration : \rdemo{numero}{titre}{contenu (items)}{fichier(s) a ouvrir} ----
\newcommand{\rdemo}[4]{%
\begin{frame}{D\'emo #1 --- #2}
  \begin{tikzpicture}[remember picture, overlay]
    \node[anchor=north east, fill=qaaccent2, text=qabg, font=\scriptsize\bfseries,
      rounded corners=3pt, inner sep=4pt] at ([xshift=-0.4cm, yshift=-0.35cm]current page.north east) {D\'EMO};
  \end{tikzpicture}
  \begin{columns}[T]
    \begin{column}{0.62\textwidth}
      \begin{itemize}
        #3
      \end{itemize}
    \end{column}
    \begin{column}{0.34\textwidth}
      \begin{datacard}
        \raggedright\scriptsize
        {\color{qamuted} Dans l'IDE}\\[1mm]
        #4
      \end{datacard}
    \end{column}
  \end{columns}
\end{frame}}

% ---- Diapo de fin « Questions ? » ----
\newcommand{\rcontact}{chouki.tibermacine@univ-ubs.fr}
\newcommand{\rquestionslide}{%
\begin{frame}{Questions ?}
  \begin{center}
    \vspace{2mm}
    \begin{tikzpicture}
      \node[regular polygon, regular polygon sides=6, draw=qaline, thick,
        minimum size=4.6cm, rotate=90] (hexin) {};
      \node[regular polygon, regular polygon sides=6, draw=qaaccent, dashed,
        thick, minimum size=6.2cm, rotate=90] (hexout) {};
      \node[font=\fontsize{92}{92}\selectfont\bfseries, color=qaaccent2] at (hexin) {?};
    \end{tikzpicture}\\[5mm]
    {\small\color{qamuted} \rcontact}
  \end{center}
\end{frame}}

%   \graphicspath{{../../latex/common/assets/}{img/}{figs/}}
% ---------------------------------------------------------------
\usetheme{metropolis}
\usepackage{appendixnumberbeamer}

\usepackage[T1]{fontenc}
\usepackage[utf8]{inputenc}
\usepackage[french]{babel}
\usepackage[scale=0.94]{tgheros}
\renewcommand{\familydefault}{\sfdefault}
\usepackage{amsmath,amssymb,amsfonts}
\usepackage{graphicx}
\usepackage{booktabs}
\usepackage{multirow}
\usepackage{tikz}
\usetikzlibrary{shapes.geometric, shapes, arrows, arrows.meta, positioning, calc, fit}
\usepackage{xcolor}
\usepackage{listings}
\usepackage[skins,listings]{tcolorbox}
\usepackage{qrcode}
\usepackage{eurosym}
\usepackage{hyperref}

% ---------------------------------------------------------------
% PALETTE (dark-mode tech: near-black canvas, violet (logo IUT) + gold accents)
% ---------------------------------------------------------------
\definecolor{qabg}{HTML}{0E1016}      % canvas background
\definecolor{qapanel}{HTML}{181B24}   % card / block panel
\definecolor{qapanelalt}{HTML}{20242F}% secondary panel tint
\definecolor{qaline}{HTML}{2B303D}    % hairline rules / table rules
\definecolor{qatext}{HTML}{ECEEF2}    % primary text (off-white)
\definecolor{qamuted}{HTML}{8B93A6}   % secondary / caption text
\definecolor{qaaccent}{HTML}{D93ED7}  % primary accent -- violet du logo IUT, eclairci pour fond sombre
\definecolor{qaaccent2}{HTML}{F2B84B} % secondary accent -- or/ambre
\definecolor{qacodekw}{HTML}{C792EA}  % code : mots-cles
\definecolor{qacodestr}{HTML}{C3E88D} % code : chaines
\definecolor{qacodecom}{HTML}{7F8AA3} % code : commentaires
\definecolor{qacodeann}{HTML}{F2B84B} % code : annotations / javadoc

\setbeamercolor{background canvas}{bg=qabg}
\setbeamercolor{normal text}{fg=qatext}
\setbeamercolor{palette primary}{bg=qabg, fg=qamuted}
\setbeamercolor{frametitle}{bg=qabg, fg=qatext}
\setbeamercolor{progress bar}{fg=qaaccent, bg=qaline}
\setbeamercolor{alerted text}{fg=qaaccent2}
\setbeamercolor{block title}{bg=qapanel, fg=qaaccent}
\setbeamercolor{block body}{bg=qapanel, fg=qatext}
\setbeamercolor{block title alerted}{bg=qapanel, fg=qaaccent2}
\setbeamercolor{block body alerted}{bg=qapanel, fg=qatext}
\setbeamercolor{block title example}{bg=qapanelalt, fg=qaaccent}
\setbeamercolor{block body example}{bg=qapanelalt, fg=qatext}
\setbeamercolor{example text}{fg=qaaccent}
\setbeamercolor{itemize item}{fg=qaaccent}
\setbeamercolor{itemize subitem}{fg=qaaccent2}
\setbeamercolor{enumerate item}{fg=qaaccent}
\setbeamercolor{enumerate subitem}{fg=qaaccent2}
\setbeamercolor{caption}{fg=qamuted}
\setbeamercolor{page number in head/foot}{fg=qamuted, bg=qabg}
\setbeamercolor{title}{fg=qatext}
\setbeamercolor{subtitle}{fg=qamuted}
\setbeamercolor{author}{fg=qatext}
\setbeamercolor{institute}{fg=qamuted}
\setbeamercolor{date}{fg=qamuted}
\setbeamercolor{section in toc}{fg=qaaccent}
\setbeamercolor{subsection in toc}{fg=qamuted}
% Sommaire de debut de section : section en cours en violet et gras, les autres en gris
\setbeamercolor{section in toc shaded}{fg=qamuted}
\setbeamercolor{subsection in toc shaded}{fg=qamuted}

\setbeamertemplate{blocks}[rounded][shadow=false]
\setbeamertemplate{itemize item}{\raisebox{0.15em}{\scriptsize\textcolor{qaaccent}{\ensuremath{\blacktriangleright}}}}
\setbeamertemplate{itemize subitem}{\raisebox{0.1em}{\tiny\textcolor{qaaccent2}{\ensuremath{\bullet}}}}

\setbeamerfont{frametitle}{series=\bfseries, size=\Large}
\setbeamerfont{title}{series=\bfseries, size=\Large}
\setbeamerfont{subtitle}{size=\normalsize}
\setbeamerfont{institute}{size=\scriptsize}
\setbeamerfont{date}{size=\scriptsize}
\setbeamerfont{author}{size=\small}

\metroset{sectionpage=none, subsectionpage=none, numbering=fraction, progressbar=frametitle}
\setbeamertemplate{section in toc}[sections numbered]

% ---------------------------------------------------------------
% LISTINGS -- theme sombre (lstlisting reste utilisable tel quel)
% \lstcode : police + couleur de base, a reutiliser dans les
%            surcharges locales (ex. basicstyle=\lstcode\scriptsize)
% ---------------------------------------------------------------
\newcommand{\lstcode}{\ttfamily\color{qatext}}

\lstdefinelanguage{TypeScript}{
  keywords={typeof, new, true, false, catch, function, return, null, switch, var,
    if, in, while, do, else, case, break, const, let, interface, class, extends,
    implements, enum, export, import, this, super, constructor, public, private,
    protected, static, async, await, throw, throws, try, finally, void},
  ndkeywords={boolean, any, number, string, int, long, double},
  sensitive=true,
  comment=[l]{//},
  morecomment=[s]{/*}{*/},
  morestring=[b]',
  morestring=[b]"
}

\lstset{
  language=Java,
  basicstyle=\lstcode\footnotesize,
  keywordstyle=\color{qacodekw}\bfseries,
  ndkeywordstyle=\color{qaaccent2}\bfseries,
  identifierstyle=\color{qatext},
  stringstyle=\color{qacodestr},
  commentstyle=\color{qacodecom}\itshape,
  morecomment=[s][\color{qacodeann}]{/**}{*/},
  backgroundcolor=\color{qapanelalt},
  rulecolor=\color{qaline},
  frame=single,
  framesep=3pt,
  xleftmargin=4pt,
  xrightmargin=4pt,
  keepspaces=true,
  columns=fullflexible,
  showstringspaces=false,
  breaklines=true,
  tabsize=2,
  escapeinside={(*@}{@*)},
  % accents UTF-8 dans le code (listings ne gere pas l'UTF-8 multi-octets)
  literate={é}{{\'e}}1 {è}{{\`e}}1 {ê}{{\^e}}1 {à}{{\`a}}1 {â}{{\^a}}1
           {ù}{{\`u}}1 {û}{{\^u}}1 {ô}{{\^o}}1 {î}{{\^i}}1 {ç}{{\c c}}1
           {É}{{\'E}}1 {À}{{\`A}}1 {€}{{\euro{}}}1 {…}{{\ldots}}1,
}

% ---------------------------------------------------------------
% CARD-STYLE BOXES
% ---------------------------------------------------------------
\tcbset{
  figcard/.style={colback=white, colframe=white, boxrule=0pt, arc=4pt,
    left=3pt, right=3pt, top=3pt, bottom=3pt, boxsep=0pt,
    halign=center, valign=center},
  datacard/.style={colback=qapanel, colframe=qaline, boxrule=0.4pt, arc=3pt,
    left=8pt, right=8pt, top=6pt, bottom=6pt, coltext=qatext,
    before upper={\color{qatext}}, halign=center},
  codebox/.style={colback=qapanelalt, colframe=qaline, boxrule=0.4pt, arc=3pt,
    left=8pt, right=8pt, top=5pt, bottom=5pt, listing only,
    listing options={basicstyle=\ttfamily\small\color{qaaccent},
      breaklines=true, showstringspaces=false, columns=fullflexible,
      keepspaces=true, frame=none, backgroundcolor={}}}
}
\newtcolorbox{figcard}[1][]{figcard,#1}
\newtcolorbox{datacard}{datacard}
\newtcblisting{codebox}{codebox}

% ---------------------------------------------------------------
% HANDY MACROS
% ---------------------------------------------------------------
% Mention d'elaboration (discrete : plus sombre que le texte secondaire)
\definecolor{qaaimention}{HTML}{4E5566}
\newcommand{\raimention}{{\tiny\color{qaaimention}Supports \'elabor\'es avec l'aide de Claude Code Opus 5.5}}
\newcommand{\hl}[1]{\alert{#1}}
\newcommand{\source}[1]{{\scriptsize\color{qamuted}#1}}

% Sommaire automatique en debut de section
\AtBeginSection[]{%
  \begin{frame}<beamer>{Plan du cours}
    \tableofcontents[currentsection]
  \end{frame}}

% Page de titre : \rtitleslide{numero du cours}{titre du cours}
% Les auteurs sont repris de \rauthors (modifiable avant \begin{document})
\newcommand{\rauthors}{Chouki Tibermacine}
\newcommand{\rtitleslide}[2]{%
\begin{frame}[plain,noframenumbering]
  \begin{center}
    \vspace{2mm}
    {\color{qaaccent}\rule{2.1cm}{1.2pt}}\\[6mm]
    {\Large\bfseries\color{qatext} R4.02}\\[1mm]
    {\Large\bfseries\color{qatext} Qualit\'e de d\'eveloppement}\\[5mm]
    {\large\color{qamuted} Cours #1}\\[1mm]
    {\large\color{qaaccent} #2}\\[6mm]
    \includegraphics[height=2.1cm]{iut_logo_transparent.png}

    \vspace{4mm}

    {\small\color{qatext} \rauthors}\\[1mm]
    {\scriptsize\color{qamuted} BUT 2 Informatique --- IUT de Vannes --- Universit\'e Bretagne Sud}\\[2mm]
    \raimention
  \end{center}
\end{frame}
}

% ---- Diapo de demonstration : \rdemo{numero}{titre}{contenu (items)}{fichier(s) a ouvrir} ----
\newcommand{\rdemo}[4]{%
\begin{frame}{D\'emo #1 --- #2}
  \begin{tikzpicture}[remember picture, overlay]
    \node[anchor=north east, fill=qaaccent2, text=qabg, font=\scriptsize\bfseries,
      rounded corners=3pt, inner sep=4pt] at ([xshift=-0.4cm, yshift=-0.35cm]current page.north east) {D\'EMO};
  \end{tikzpicture}
  \begin{columns}[T]
    \begin{column}{0.62\textwidth}
      \begin{itemize}
        #3
      \end{itemize}
    \end{column}
    \begin{column}{0.34\textwidth}
      \begin{datacard}
        \raggedright\scriptsize
        {\color{qamuted} Dans l'IDE}\\[1mm]
        #4
      \end{datacard}
    \end{column}
  \end{columns}
\end{frame}}

% ---- Diapo de fin « Questions ? » ----
\newcommand{\rcontact}{chouki.tibermacine@univ-ubs.fr}
\newcommand{\rquestionslide}{%
\begin{frame}{Questions ?}
  \begin{center}
    \vspace{2mm}
    \begin{tikzpicture}
      \node[regular polygon, regular polygon sides=6, draw=qaline, thick,
        minimum size=4.6cm, rotate=90] (hexin) {};
      \node[regular polygon, regular polygon sides=6, draw=qaaccent, dashed,
        thick, minimum size=6.2cm, rotate=90] (hexout) {};
      \node[font=\fontsize{92}{92}\selectfont\bfseries, color=qaaccent2] at (hexin) {?};
    \end{tikzpicture}\\[5mm]
    {\small\color{qamuted} \rcontact}
  \end{center}
\end{frame}}

%   \graphicspath{{../../latex/common/assets/}{img/}{figs/}}
% ---------------------------------------------------------------
\usetheme{metropolis}
\usepackage{appendixnumberbeamer}

\usepackage[T1]{fontenc}
\usepackage[utf8]{inputenc}
\usepackage[french]{babel}
\usepackage[scale=0.94]{tgheros}
\renewcommand{\familydefault}{\sfdefault}
\usepackage{amsmath,amssymb,amsfonts}
\usepackage{graphicx}
\usepackage{booktabs}
\usepackage{multirow}
\usepackage{tikz}
\usetikzlibrary{shapes.geometric, shapes, arrows, arrows.meta, positioning, calc, fit}
\usepackage{xcolor}
\usepackage{listings}
\usepackage[skins,listings]{tcolorbox}
\usepackage{qrcode}
\usepackage{eurosym}
\usepackage{hyperref}

% ---------------------------------------------------------------
% PALETTE (dark-mode tech: near-black canvas, violet (logo IUT) + gold accents)
% ---------------------------------------------------------------
\definecolor{qabg}{HTML}{0E1016}      % canvas background
\definecolor{qapanel}{HTML}{181B24}   % card / block panel
\definecolor{qapanelalt}{HTML}{20242F}% secondary panel tint
\definecolor{qaline}{HTML}{2B303D}    % hairline rules / table rules
\definecolor{qatext}{HTML}{ECEEF2}    % primary text (off-white)
\definecolor{qamuted}{HTML}{8B93A6}   % secondary / caption text
\definecolor{qaaccent}{HTML}{D93ED7}  % primary accent -- violet du logo IUT, eclairci pour fond sombre
\definecolor{qaaccent2}{HTML}{F2B84B} % secondary accent -- or/ambre
\definecolor{qacodekw}{HTML}{C792EA}  % code : mots-cles
\definecolor{qacodestr}{HTML}{C3E88D} % code : chaines
\definecolor{qacodecom}{HTML}{7F8AA3} % code : commentaires
\definecolor{qacodeann}{HTML}{F2B84B} % code : annotations / javadoc

\setbeamercolor{background canvas}{bg=qabg}
\setbeamercolor{normal text}{fg=qatext}
\setbeamercolor{palette primary}{bg=qabg, fg=qamuted}
\setbeamercolor{frametitle}{bg=qabg, fg=qatext}
\setbeamercolor{progress bar}{fg=qaaccent, bg=qaline}
\setbeamercolor{alerted text}{fg=qaaccent2}
\setbeamercolor{block title}{bg=qapanel, fg=qaaccent}
\setbeamercolor{block body}{bg=qapanel, fg=qatext}
\setbeamercolor{block title alerted}{bg=qapanel, fg=qaaccent2}
\setbeamercolor{block body alerted}{bg=qapanel, fg=qatext}
\setbeamercolor{block title example}{bg=qapanelalt, fg=qaaccent}
\setbeamercolor{block body example}{bg=qapanelalt, fg=qatext}
\setbeamercolor{example text}{fg=qaaccent}
\setbeamercolor{itemize item}{fg=qaaccent}
\setbeamercolor{itemize subitem}{fg=qaaccent2}
\setbeamercolor{enumerate item}{fg=qaaccent}
\setbeamercolor{enumerate subitem}{fg=qaaccent2}
\setbeamercolor{caption}{fg=qamuted}
\setbeamercolor{page number in head/foot}{fg=qamuted, bg=qabg}
\setbeamercolor{title}{fg=qatext}
\setbeamercolor{subtitle}{fg=qamuted}
\setbeamercolor{author}{fg=qatext}
\setbeamercolor{institute}{fg=qamuted}
\setbeamercolor{date}{fg=qamuted}
\setbeamercolor{section in toc}{fg=qaaccent}
\setbeamercolor{subsection in toc}{fg=qamuted}
% Sommaire de debut de section : section en cours en violet et gras, les autres en gris
\setbeamercolor{section in toc shaded}{fg=qamuted}
\setbeamercolor{subsection in toc shaded}{fg=qamuted}

\setbeamertemplate{blocks}[rounded][shadow=false]
\setbeamertemplate{itemize item}{\raisebox{0.15em}{\scriptsize\textcolor{qaaccent}{\ensuremath{\blacktriangleright}}}}
\setbeamertemplate{itemize subitem}{\raisebox{0.1em}{\tiny\textcolor{qaaccent2}{\ensuremath{\bullet}}}}

\setbeamerfont{frametitle}{series=\bfseries, size=\Large}
\setbeamerfont{title}{series=\bfseries, size=\Large}
\setbeamerfont{subtitle}{size=\normalsize}
\setbeamerfont{institute}{size=\scriptsize}
\setbeamerfont{date}{size=\scriptsize}
\setbeamerfont{author}{size=\small}

\metroset{sectionpage=none, subsectionpage=none, numbering=fraction, progressbar=frametitle}
\setbeamertemplate{section in toc}[sections numbered]

% ---------------------------------------------------------------
% LISTINGS -- theme sombre (lstlisting reste utilisable tel quel)
% \lstcode : police + couleur de base, a reutiliser dans les
%            surcharges locales (ex. basicstyle=\lstcode\scriptsize)
% ---------------------------------------------------------------
\newcommand{\lstcode}{\ttfamily\color{qatext}}

\lstdefinelanguage{TypeScript}{
  keywords={typeof, new, true, false, catch, function, return, null, switch, var,
    if, in, while, do, else, case, break, const, let, interface, class, extends,
    implements, enum, export, import, this, super, constructor, public, private,
    protected, static, async, await, throw, throws, try, finally, void},
  ndkeywords={boolean, any, number, string, int, long, double},
  sensitive=true,
  comment=[l]{//},
  morecomment=[s]{/*}{*/},
  morestring=[b]',
  morestring=[b]"
}

\lstset{
  language=Java,
  basicstyle=\lstcode\footnotesize,
  keywordstyle=\color{qacodekw}\bfseries,
  ndkeywordstyle=\color{qaaccent2}\bfseries,
  identifierstyle=\color{qatext},
  stringstyle=\color{qacodestr},
  commentstyle=\color{qacodecom}\itshape,
  morecomment=[s][\color{qacodeann}]{/**}{*/},
  backgroundcolor=\color{qapanelalt},
  rulecolor=\color{qaline},
  frame=single,
  framesep=3pt,
  xleftmargin=4pt,
  xrightmargin=4pt,
  keepspaces=true,
  columns=fullflexible,
  showstringspaces=false,
  breaklines=true,
  tabsize=2,
  escapeinside={(*@}{@*)},
  % accents UTF-8 dans le code (listings ne gere pas l'UTF-8 multi-octets)
  literate={é}{{\'e}}1 {è}{{\`e}}1 {ê}{{\^e}}1 {à}{{\`a}}1 {â}{{\^a}}1
           {ù}{{\`u}}1 {û}{{\^u}}1 {ô}{{\^o}}1 {î}{{\^i}}1 {ç}{{\c c}}1
           {É}{{\'E}}1 {À}{{\`A}}1 {€}{{\euro{}}}1 {…}{{\ldots}}1,
}

% ---------------------------------------------------------------
% CARD-STYLE BOXES
% ---------------------------------------------------------------
\tcbset{
  figcard/.style={colback=white, colframe=white, boxrule=0pt, arc=4pt,
    left=3pt, right=3pt, top=3pt, bottom=3pt, boxsep=0pt,
    halign=center, valign=center},
  datacard/.style={colback=qapanel, colframe=qaline, boxrule=0.4pt, arc=3pt,
    left=8pt, right=8pt, top=6pt, bottom=6pt, coltext=qatext,
    before upper={\color{qatext}}, halign=center},
  codebox/.style={colback=qapanelalt, colframe=qaline, boxrule=0.4pt, arc=3pt,
    left=8pt, right=8pt, top=5pt, bottom=5pt, listing only,
    listing options={basicstyle=\ttfamily\small\color{qaaccent},
      breaklines=true, showstringspaces=false, columns=fullflexible,
      keepspaces=true, frame=none, backgroundcolor={}}}
}
\newtcolorbox{figcard}[1][]{figcard,#1}
\newtcolorbox{datacard}{datacard}
\newtcblisting{codebox}{codebox}

% ---------------------------------------------------------------
% HANDY MACROS
% ---------------------------------------------------------------
% Mention d'elaboration (discrete : plus sombre que le texte secondaire)
\definecolor{qaaimention}{HTML}{4E5566}
\newcommand{\raimention}{{\tiny\color{qaaimention}Supports \'elabor\'es avec l'aide de Claude Code Opus 5.5}}
\newcommand{\hl}[1]{\alert{#1}}
\newcommand{\source}[1]{{\scriptsize\color{qamuted}#1}}

% Sommaire automatique en debut de section
\AtBeginSection[]{%
  \begin{frame}<beamer>{Plan du cours}
    \tableofcontents[currentsection]
  \end{frame}}

% Page de titre : \rtitleslide{numero du cours}{titre du cours}
% Les auteurs sont repris de \rauthors (modifiable avant \begin{document})
\newcommand{\rauthors}{Chouki Tibermacine}
\newcommand{\rtitleslide}[2]{%
\begin{frame}[plain,noframenumbering]
  \begin{center}
    \vspace{2mm}
    {\color{qaaccent}\rule{2.1cm}{1.2pt}}\\[6mm]
    {\Large\bfseries\color{qatext} R4.02}\\[1mm]
    {\Large\bfseries\color{qatext} Qualit\'e de d\'eveloppement}\\[5mm]
    {\large\color{qamuted} Cours #1}\\[1mm]
    {\large\color{qaaccent} #2}\\[6mm]
    \includegraphics[height=2.1cm]{iut_logo_transparent.png}

    \vspace{4mm}

    {\small\color{qatext} \rauthors}\\[1mm]
    {\scriptsize\color{qamuted} BUT 2 Informatique --- IUT de Vannes --- Universit\'e Bretagne Sud}\\[2mm]
    \raimention
  \end{center}
\end{frame}
}

% ---- Diapo de demonstration : \rdemo{numero}{titre}{contenu (items)}{fichier(s) a ouvrir} ----
\newcommand{\rdemo}[4]{%
\begin{frame}{D\'emo #1 --- #2}
  \begin{tikzpicture}[remember picture, overlay]
    \node[anchor=north east, fill=qaaccent2, text=qabg, font=\scriptsize\bfseries,
      rounded corners=3pt, inner sep=4pt] at ([xshift=-0.4cm, yshift=-0.35cm]current page.north east) {D\'EMO};
  \end{tikzpicture}
  \begin{columns}[T]
    \begin{column}{0.62\textwidth}
      \begin{itemize}
        #3
      \end{itemize}
    \end{column}
    \begin{column}{0.34\textwidth}
      \begin{datacard}
        \raggedright\scriptsize
        {\color{qamuted} Dans l'IDE}\\[1mm]
        #4
      \end{datacard}
    \end{column}
  \end{columns}
\end{frame}}

% ---- Diapo de fin « Questions ? » ----
\newcommand{\rcontact}{chouki.tibermacine@univ-ubs.fr}
\newcommand{\rquestionslide}{%
\begin{frame}{Questions ?}
  \begin{center}
    \vspace{2mm}
    \begin{tikzpicture}
      \node[regular polygon, regular polygon sides=6, draw=qaline, thick,
        minimum size=4.6cm, rotate=90] (hexin) {};
      \node[regular polygon, regular polygon sides=6, draw=qaaccent, dashed,
        thick, minimum size=6.2cm, rotate=90] (hexout) {};
      \node[font=\fontsize{92}{92}\selectfont\bfseries, color=qaaccent2] at (hexin) {?};
    \end{tikzpicture}\\[5mm]
    {\small\color{qamuted} \rcontact}
  \end{center}
\end{frame}}

\graphicspath{{../../latex/common/assets/}{img/}}

\title{Testabilité et doublures de test}
\author[C. Tibermacine]{Chouki Tibermacine}
\date{}

\lstdefinestyle{sortie}{language={}, basicstyle=\lstcode\scriptsize, keywordstyle={}, commentstyle={}, stringstyle={}}

\tikzset{
  comp/.style={rectangle, rounded corners=2pt, draw=qaaccent, thick, fill=qapanel, text=qatext,
    minimum height=0.6cm, minimum width=2.6cm, align=center, font=\scriptsize},
  dep/.style={rectangle, rounded corners=2pt, draw=qamuted, thick, fill=qapanelalt, text=qatext,
    minimum height=0.55cm, minimum width=2.4cm, align=center, font=\scriptsize},
  dbl/.style={rectangle, rounded corners=2pt, draw=qaaccent2, thick, dashed, fill=qapanelalt, text=qatext,
    minimum height=0.55cm, minimum width=2.4cm, align=center, font=\scriptsize},
  arr/.style={-Latex, qamuted, thick}
}

\begin{document}

\rtitleslide{3}{Testabilité et doublures de test}

\begin{frame}{Le problème}
  \begin{columns}[c]
    \begin{column}{0.5\textwidth}
      \begin{center}
      \begin{tikzpicture}
        \node[comp] (svc) at (0,0) {\texttt{VirementService}};
        \node[dep] (repo)  at (-1.75, 1.6)  {\texttt{CompteRepository}\\{\tiny\color{qamuted}PostgreSQL}};
        \node[dep] (taux)  at ( 1.75, 1.6)  {\texttt{TauxChangeClient}\\{\tiny\color{qamuted}API HTTP externe}};
        \node[dep] (notif) at (-1.75,-1.6)  {\texttt{NotificationService}\\{\tiny\color{qamuted}envoie de vrais SMS}};
        \node[dep] (clock) at ( 1.75,-1.6)  {\texttt{Clock}\\{\tiny\color{qamuted}la date change}};
        \draw[arr] ([xshift=-0.6cm]svc.north) -- (repo.south);
        \draw[arr] ([xshift=0.6cm]svc.north)  -- (taux.south);
        \draw[arr] ([xshift=-0.6cm]svc.south) -- (notif.north);
        \draw[arr] ([xshift=0.6cm]svc.south)  -- (clock.north);
      \end{tikzpicture}
      \end{center}
    \end{column}
    \begin{column}{0.48\textwidth}
      On veut tester \texttt{VirementService.virer} \textbf{unitairement}. Mais avec ses vraies dépendances, le test :
      \begin{itemize}
        \item a besoin d'une base démarrée et remplie \alert{(lent, pas F)}
        \item dépend du taux du jour et du réseau \alert{(pas R)}
        \item envoie un vrai SMS à chaque exécution
        \item ne peut pas provoquer une panne de l'API
        \item donne un résultat différent le dimanche
      \end{itemize}
      \vspace{1mm}
      $\Rightarrow$ remplacer les dépendances par des \textbf{doublures}, comme au cinéma
    \end{column}
  \end{columns}
\end{frame}

% ================================================================
\section{Écrire du code testable}

\begin{frame}[fragile]{Un code difficile à tester}
\begin{lstlisting}[basicstyle=\lstcode\scriptsize]
public class VirementService {
    private final CompteRepository comptes = new JpaCompteRepository();        // vraie base
    private final TauxChangeClient taux = new HttpTauxChangeClient("https://...");

    public void virer(String numSource, String numCible, BigDecimal montant) {
        if (LocalDate.now().getDayOfWeek() == DayOfWeek.SUNDAY) {               // date réelle
            throw new VirementImpossibleException("pas de virement le dimanche");
        }
        // ... débit, crédit, conversion avec taux.taux(...)
        new SmsNotifier().envoyer(numSource, "Virement effectué");              // vrai SMS
    }
}
\end{lstlisting}
  \begin{itemize}
    \item Chaque \texttt{new} et chaque \texttt{LocalDate.now()} est une dépendance \textbf{cachée} et \textbf{figée}
    \item Le test n'a aucun moyen de les remplacer : il subit la base, le réseau et le calendrier
    \item Ce n'est pas un problème de test : c'est un problème de \textbf{conception} (couplage fort)
  \end{itemize}
\end{frame}

\rdemo{1}{Le code en dur ne se teste pas}{%
  \item On lance un test de \texttt{VirementServiceLegacy}\ldots
  \item \ldots{} il échoue : \textit{connexion refusée à PostgreSQL}
  \item L'échec n'a \textbf{rien à voir} avec la logique du virement
  \item Et si la base était démarrée ? API de change injoignable, \textbf{vrai SMS} envoyé, échec le dimanche
}{\texttt{legacy/VirementServiceLegacy}\\[1mm]\texttt{demo1/VirementServiceLegacyTest}}

\begin{frame}[fragile]{Le même code, testable}
\begin{lstlisting}[basicstyle=\lstcode\scriptsize]
public class VirementService {
    private final CompteRepository comptes;
    private final TauxChangeClient taux;
    private final NotificationService notifications;
    private final Clock horloge;
    public VirementService(CompteRepository comptes, TauxChangeClient taux,
                           NotificationService notifications, Clock horloge) {
        this.comptes = comptes;  this.taux = taux;
        this.notifications = notifications;  this.horloge = horloge;
    }
    public void virer(String numSource, String numCible, BigDecimal montant) {
        if (LocalDate.now(horloge).getDayOfWeek() == SUNDAY) { ... }
        notifications.notifier(numSource, "Virement effectué");
    }
}
\end{lstlisting}
  \small
  \textbf{Injection par constructeur}, dépendances typées par des \textbf{interfaces}, horloge en paramètre (\texttt{Clock}).
  En production, Spring injecte les vraies implémentations ; en test, \textbf{vous} injectez des doublures.
\end{frame}

\begin{frame}{Vous l'avez déjà fait : le patron Repository}
  \begin{columns}[c]
    \begin{column}{0.5\textwidth}
      \begin{center}
      \begin{tikzpicture}[node distance=0.5cm]
        \node[comp] (svc) {\texttt{VirementService}};
        \node[comp, below=of svc, draw=qaaccent2] (itf) {\guillemotleft interface\guillemotright\\\texttt{CompteRepository}};
        \node[dep, below left=0.7cm and -0.8cm of itf] (jpa) {implémentation\\Spring Data / JPA};
        \node[dbl, below right=0.7cm and -0.8cm of itf] (mem) {implémentation\\en mémoire (test)};
        \draw[arr] (svc) -- node[right, font=\tiny, text=qamuted] {dépend de} (itf);
        \draw[arr, dashed] (jpa) -- (itf);
        \draw[arr, dashed] (mem) -- (itf);
      \end{tikzpicture}
      \end{center}
    \end{column}
    \begin{column}{0.48\textwidth}
      \begin{itemize}
        \item En R3.07, le Repository cachait JPA derrière une interface « métier »
        \item Le service ne connaît que l'\textbf{interface} : il ne sait pas s'il parle à PostgreSQL ou à une \texttt{HashMap}
        \item C'est exactement ce qui le rend testable
        \item Règle générale : dépendre d'\textbf{abstractions}, pas de classes concrètes (principe d'inversion des dépendances)
      \end{itemize}
    \end{column}
  \end{columns}
\end{frame}

% ================================================================
\section{Les doublures de test}

\begin{frame}{Cinq sortes de doublures (Meszaros, 2007)}
  \footnotesize
  \begin{datacard}
    \begin{tabular}{@{}p{1.2cm}p{4.9cm}p{6.0cm}@{}}
      \toprule
      \textbf{Nom} & \textbf{Rôle} & \textbf{Exemple dans la banque} \\
      \midrule
      \textbf{\color{qaaccent}Dummy} & passé en paramètre mais jamais utilisé & un \texttt{NotificationService} \texttt{null} dans un test où aucune notification ne part \\[1mm]
      \textbf{\color{qaaccent}Stub} & renvoie des réponses \textbf{préparées} & un \texttt{TauxChangeClient} qui renvoie toujours 1,10 \\[1mm]
      \textbf{\color{qaaccent}Fake} & implémentation \textbf{simplifiée} mais qui marche & un \texttt{CompteRepository} dans une \texttt{HashMap} (le \texttt{ContexteBanque} du TP2) \\[1mm]
      \textbf{\color{qaaccent}Spy} & \textbf{enregistre} les appels reçus, pour les vérifier après & un \texttt{NotificationService} qui garde la liste des messages « envoyés » \\[1mm]
      \textbf{\color{qaaccent}Mock} & programmé avec des \textbf{attentes} sur les appels ; le test échoue si elles ne sont pas respectées & « \texttt{notifier} doit être appelé une fois avec \texttt{FR76-001} » \\
      \bottomrule
    \end{tabular}
  \end{datacard}
  \vspace{1mm}
  \source{Gerard Meszaros, \textit{xUnit Test Patterns}, 2007. Dans le langage courant, on dit souvent « mock » pour toutes les doublures.}
\end{frame}

\begin{frame}[fragile]{Des doublures écrites à la main}
\begin{lstlisting}[basicstyle=\lstcode\scriptsize]
// Stub : toujours la même réponse (TauxChangeClient n'a qu'une méthode : une lambda suffit)
TauxChangeClient taux = (de, vers) -> new BigDecimal("1.10");

// Fake : une vraie logique, mais sans base de données
class CompteRepositoryEnMemoire implements CompteRepository {
    private final Map<String, Compte> comptes = new HashMap<>();

    public Optional<Compte> findByNumero(String numero) {
        return Optional.ofNullable(comptes.get(numero));
    }
    public Compte save(Compte compte) {
        comptes.put(compte.getNumero(), compte);
        return compte;
    }
}
\end{lstlisting}
  \small
  \begin{itemize}
    \item Un spy à la main : une implémentation de \texttt{NotificationService} qui ajoute chaque message reçu dans une liste, qu'on vérifie à la fin du test
    \item Simple pour de petites interfaces ; fastidieux pour 10 méthodes et 30 tests : on utilise une bibliothèque
  \end{itemize}
\end{frame}

\rdemo{2}{Stub, fake et spy faits main}{%
  \item Un fake (\texttt{HashMap}), un spy (liste de messages), un stub (une lambda)
  \item 4 tests verts en quelques millisecondes, sans base ni réseau
  \item Tester « le dimanche » un lundi matin, grâce à \texttt{Clock}
  \item Piège du fuseau : \texttt{2027-03-06T23:30Z}, samedi en UTC\ldots{} dimanche à Paris
}{\texttt{demo2/CompteRepositoryEnMemoire}\\\texttt{demo2/NotificationsEspion}\\\texttt{demo2/DoubluresALaMainTest}}

\begin{frame}{Vérifier un état ou un comportement ?}
  \small
  \begin{columns}[T]
    \begin{column}{0.48\textwidth}
      \begin{block}{Vérification d'état}
        \begin{itemize}
          \item Après l'action, on regarde le \textbf{résultat} : soldes, valeur renvoyée
          \item Utilise des vrais objets, des stubs et des fakes
          \item Résiste bien aux changements d'implémentation
          \item École « classique » (Detroit ou Chicago, Kent Beck)
        \end{itemize}
      \end{block}
    \end{column}
    \begin{column}{0.48\textwidth}
      \begin{block}{Vérification de comportement}
        \begin{itemize}
          \item On vérifie que les bons \textbf{appels} ont eu lieu, avec les bons paramètres
          \item Utilise des mocks et des spies
          \item Indispensable quand l'effet est \textbf{à l'extérieur} : un SMS envoyé n'a pas d'« état » à lire
          \item École « mockiste » (Londres, Freeman et Pryce)
        \end{itemize}
      \end{block}
    \end{column}
  \end{columns}
  \vspace{3mm}
  \begin{exampleblock}{Une règle simple}
    \textbf{Simuler les requêtes} (ce qui renvoie une valeur : \texttt{findByNumero}, \texttt{taux}) avec des stubs ; \textbf{vérifier les commandes} (ce qui produit un effet : \texttt{notifier}, \texttt{save}).
  \end{exampleblock}
\end{frame}

\begin{frame}{Quelle doublure choisir ?}
  \small
  \begin{datacard}
    \begin{tabular}{@{}p{6.4cm}p{6.2cm}@{}}
      \toprule
      \textbf{Situation} & \textbf{Doublure} \\
      \midrule
      Le paramètre est obligatoire mais ne sert pas dans ce test & \textbf{dummy} (souvent \texttt{null} ou un mock jamais programmé) \\[0.5mm]
      J'ai besoin qu'une requête renvoie une valeur précise, ou échoue & \textbf{stub} (\texttt{when ... thenReturn / thenThrow}) \\[0.5mm]
      Le collaborateur a un vrai comportement (stocker puis relire) utilisé par plusieurs tests & \textbf{fake} (en mémoire) \\[0.5mm]
      Je dois vérifier qu'un effet extérieur a eu lieu (SMS, e-mail, enregistrement) & \textbf{mock} ou \textbf{spy} (\texttt{verify}) \\[0.5mm]
      Le collaborateur est un objet du domaine (\texttt{Compte}) ou une valeur (\texttt{BigDecimal}) & \alert{aucune} : on utilise le vrai \\[0.5mm]
      Je dois vérifier que la vraie base ou la vraie API marche & \alert{aucune} : test d'\textbf{intégration} (CM4) \\
      \bottomrule
    \end{tabular}
  \end{datacard}
\end{frame}

% ================================================================
\section{Mockito}

\begin{frame}[fragile]{Mettre en place Mockito}
  \begin{columns}[T]
    \begin{column}{0.6\textwidth}
\begin{lstlisting}[basicstyle=\lstcode\scriptsize]
@ExtendWith(MockitoExtension.class)
class VirementServiceTest {

    @Mock CompteRepository comptes;
    @Mock TauxChangeClient taux;
    @Mock NotificationService notifications;

    Clock horloge = Clock.fixed(
        Instant.parse("2027-03-01T10:00:00Z"), ZoneOffset.UTC);

    VirementService service;

    @BeforeEach
    void init() {
        service = new VirementService(comptes, taux,
                                      notifications, horloge);
    }
}
\end{lstlisting}
    \end{column}
    \begin{column}{0.37\textwidth}
      \small
      \begin{itemize}
        \item \texttt{@Mock} : Mockito crée une implémentation de l'interface
        \item Par défaut, ses méthodes ne font rien et renvoient une valeur « vide » (\texttt{null}, 0, \texttt{Optional.empty()}, liste vide)
        \item \texttt{Clock.fixed} : toujours le lundi 1\textsuperscript{er} mars 2027
        \item \texttt{@InjectMocks} existe, mais n'injecte que des mocks (pas la \texttt{Clock}) : construire le service soi-même est plus clair
      \end{itemize}
    \end{column}
  \end{columns}
\end{frame}

\begin{frame}[fragile]{Préparer les réponses : \texttt{when ... thenReturn}}
\begin{lstlisting}[basicstyle=\lstcode\scriptsize]
@Test
void virer_comptesEnEurosEtEnDollars_convertitLeMontant() {
    Compte source = new Compte("FR76-001", Devise.EUR, eur("500.00"));
    Compte cible  = new Compte("US-042",   Devise.USD, eur("0.00"));
    when(comptes.findByNumero("FR76-001")).thenReturn(Optional.of(source));   // Arrange : stubs
    when(comptes.findByNumero("US-042")).thenReturn(Optional.of(cible));
    when(taux.taux(Devise.EUR, Devise.USD)).thenReturn(new BigDecimal("1.10"));

    service.virer("FR76-001", "US-042", eur("100.00"));                       // Act

    assertThat(source.getSolde()).isEqualByComparingTo("400.00");             // Assert : état
    assertThat(cible.getSolde()).isEqualByComparingTo("110.00");
}
\end{lstlisting}
  \begin{itemize}
    \item On ne programme que ce dont \textbf{ce} test a besoin ; avec \texttt{MockitoExtension}, un stub inutilisé fait échouer le test (\texttt{UnnecessaryStubbingException})
    \item \texttt{thenReturn(a, b, c)} : réponses successives ; \texttt{thenAnswer(...)} : réponse calculée
  \end{itemize}
\end{frame}

\begin{frame}[fragile]{Simuler une panne : \texttt{thenThrow}}
\begin{lstlisting}[basicstyle=\lstcode\scriptsize]
@Test
void virer_apiDeChangeIndisponible_refuseSansRienModifier() {
    Compte source = new Compte("FR76-001", Devise.EUR, eur("500.00"));
    Compte cible  = new Compte("US-042",   Devise.USD, eur("0.00"));
    when(comptes.findByNumero("FR76-001")).thenReturn(Optional.of(source));
    when(comptes.findByNumero("US-042")).thenReturn(Optional.of(cible));
    when(taux.taux(any(), any())).thenThrow(new TauxIndisponibleException("délai dépassé"));

    assertThatThrownBy(() -> service.virer("FR76-001", "US-042", eur("100.00")))
        .isInstanceOf(VirementImpossibleException.class);

    assertThat(source.getSolde()).isEqualByComparingTo("500.00");   // rien n'a bougé
    verify(comptes, never()).save(any());                            // rien n'a été enregistré
    verifyNoInteractions(notifications);                              // personne n'a été prévenu
}
\end{lstlisting}
  \begin{itemize}
    \item Les pannes (réseau coupé, délai dépassé, service en maintenance) sont \textbf{impossibles à provoquer} de façon fiable avec la vraie API
    \item Avec une doublure, c'est une ligne : les cas d'erreur deviennent faciles à tester
  \end{itemize}
\end{frame}

\begin{frame}[fragile]{Vérifier les appels : \texttt{verify}}
  \begin{columns}[T]
    \begin{column}{0.55\textwidth}
\begin{lstlisting}[basicstyle=\lstcode\scriptsize]
service.virer("FR76-001", "FR76-002", eur("1500.00"));

verify(notifications).notifier(eq("FR76-001"), anyString());
verify(comptes, times(2)).save(any(Compte.class));
verify(taux, never()).taux(any(), any());  // même devise
\end{lstlisting}
\begin{lstlisting}[basicstyle=\lstcode\scriptsize]
// Capturer un argument pour l'examiner
ArgumentCaptor<String> message =
    ArgumentCaptor.forClass(String.class);
verify(notifications).notifier(eq("FR76-001"),
                               message.capture());
assertThat(message.getValue()).contains("1500.00 €");
\end{lstlisting}
    \end{column}
    \begin{column}{0.42\textwidth}
      \begin{itemize}
        \item \texttt{verify(mock)} : appelé exactement une fois avec ces arguments
        \item \texttt{times(n)}, \texttt{never()}, \texttt{atLeastOnce()}, \texttt{verifyNoInteractions}
        \item \textbf{Matchers} : \texttt{any()}, \texttt{anyString()}, \texttt{eq(...)}, \texttt{argThat(...)}
          \begin{itemize}
            \item si un argument utilise un matcher, \alert{tous} doivent en utiliser un (d'où \texttt{eq("FR76-001")})
          \end{itemize}
        \item \texttt{ArgumentCaptor} : récupérer l'argument réellement passé pour faire des assertions précises
      \end{itemize}
    \end{column}
  \end{columns}
\end{frame}

\rdemo{3}{Mockito en action}{%
  \item Que renvoie un mock qu'on n'a pas programmé ?
  \item Conversion avec un taux à 4 décimales : l'arrondi est-il bon ?
  \item Une panne d'API en une ligne (\texttt{thenThrow})
  \item \texttt{verify} et \texttt{ArgumentCaptor} : on casse le message attendu pour lire le diff
  \item Valeur limite 1\,000 \euro{} : on introduit un bug \texttt{>=}, le test le voit
}{\texttt{demo3/VirementServiceMockitoTest}\\[1mm]\texttt{service/VirementService}}

\begin{frame}[fragile]{Espionner un vrai objet : \texttt{@Spy}}
  \begin{columns}[T]
    \begin{column}{0.55\textwidth}
\begin{lstlisting}[basicstyle=\lstcode\scriptsize]
Compte source = spy(new Compte("FR76-001", eur("500.00")));
when(comptes.findByNumero("FR76-001"))
    .thenReturn(Optional.of(source));
// ...
service.virer("FR76-001", "FR76-002", eur("100.00"));

verify(source).debiter(eur("100.00"));   // vrai débit exécuté
\end{lstlisting}
    \end{column}
    \begin{column}{0.42\textwidth}
      \begin{itemize}
        \item Un \textbf{spy} Mockito enveloppe un \textbf{vrai} objet : les vraies méthodes s'exécutent, mais les appels sont enregistrés
        \item On peut aussi remplacer une seule méthode : \texttt{doReturn(x).when(spy).methode()}
        \item \alert{À utiliser rarement} : si on a besoin d'espionner un objet du domaine, c'est souvent que le test vérifie l'implémentation plutôt que le résultat
      \end{itemize}
    \end{column}
  \end{columns}
\end{frame}

\rdemo{4}{Lire les erreurs de Mockito}{%
  \item Quatre tests rouges, quatre erreurs que vous rencontrerez en TP :
    \begin{itemize}
      \item \texttt{UnnecessaryStubbingException}
      \item \texttt{InvalidUseOfMatchersException}
      \item un stub oublié\ldots{} et une exception déroutante
      \item \texttt{PotentialStubbingProblem}
    \end{itemize}
  \item Le message de Mockito contient presque toujours la solution
}{\texttt{demo4/PiegesMockitoTest}}

\begin{frame}{Les pièges du mocking}
  \footnotesize
  \begin{columns}[T]
    \begin{column}{0.48\textwidth}
      \begin{alertblock}{Tout simuler}
        \begin{itemize}
          \item Simuler \texttt{Compte}, \texttt{BigDecimal}, une \texttt{List}\ldots{} : le test ne vérifie plus que ses propres réglages
          \item On simule les \textbf{frontières} (base, réseau, horloge, e-mail), pas les objets du domaine
        \end{itemize}
      \end{alertblock}
      \begin{alertblock}{Simuler ce qu'on ne possède pas}
        \begin{itemize}
          \item Simuler le client HTTP d'une API : on reproduit \textbf{sa} documentation, peut-être de travers
          \item Mieux : l'envelopper dans \textbf{notre} interface (\texttt{TauxChangeClient}), simuler celle-ci, tester l'adaptateur réel en intégration (CM4)
        \end{itemize}
      \end{alertblock}
    \end{column}
    \begin{column}{0.48\textwidth}
      \begin{alertblock}{Tests couplés à l'implémentation}
        \begin{itemize}
          \item Un test qui vérifie chaque appel interne casse au moindre refactoring, même si le comportement est inchangé
          \item Vérifier ce qui \textbf{compte} pour l'extérieur : le résultat, et les effets visibles (notification, enregistrement)
        \end{itemize}
      \end{alertblock}
      \begin{alertblock}{Des mocks qui mentent}
        \begin{itemize}
          \item Le mock renvoie \texttt{Optional.of(...)} alors que la vraie base renverrait vide : test vert, production rouge
          \item Les doublures ne remplacent pas les \textbf{tests d'intégration} : elles les complètent
        \end{itemize}
      \end{alertblock}
    \end{column}
  \end{columns}
\end{frame}

\rdemo{5}{Un test trop couplé}{%
  \item Refactoring : deux \texttt{save} remplacés par un \texttt{saveAll}
  \item Le test qui comptait les \texttt{save} passe au rouge\ldots
  \item \ldots{} celui qui vérifie l'état enregistré (avec un fake) reste vert
  \item Le comportement a-t-il changé pour le client ? Qui avait tort ?
}{\texttt{service/VirementService}\\\texttt{demo3/\ldots\#virer\_enregistre\-LesDeuxComptes}\\\texttt{demo5/VirementServiceRobusteTest}}

\begin{frame}{À retenir}
  \small
  \begin{itemize}
    \item Un code testable reçoit ses dépendances (\textbf{injection par constructeur}) et dépend d'\textbf{interfaces} ; pas de \texttt{new} ni de \texttt{LocalDate.now()} cachés
    \item Cinq doublures : \textbf{dummy}, \textbf{stub}, \textbf{fake}, \textbf{spy}, \textbf{mock}
    \item Simuler les requêtes, vérifier les commandes ; préférer la vérification d'état quand c'est possible
    \item Mockito : \texttt{@Mock}, \texttt{when(...).thenReturn(...)}, \texttt{thenThrow}, \texttt{verify}, \texttt{never()}, \texttt{ArgumentCaptor}, \texttt{spy}
    \item Simuler les \textbf{frontières} du système, pas le domaine ; envelopper les API externes dans nos propres interfaces
    \item Les doublures rendent les tests rapides et reproductibles, et les pannes faciles à provoquer ; elles ne prouvent pas que l'intégration marche (CM4)
  \end{itemize}
  \vspace{1mm}
  \begin{exampleblock}{TP3}
    Tester \texttt{VirementService} avec Mockito, puis rendre testable une classe qui ne l'est pas. \alert{Quiz 1} en fin de semaine, lors d'une séance à part (CM1 à CM3).
  \end{exampleblock}
\end{frame}

\rquestionslide

\end{document}
